\solchapter{fluids}

\sol{fluids}{10} With no inertia, the body's velocity at each instant is a linear function of the rate of change of its shape. What does that imply about the net displacement over a path in shape space that is traced forwards and then backwards?

\sol{fluids}{11} By Kutta--Joukowski, lift per length is $\rho v\Gamma$. For a spinning ball, estimate the circulation as (surface speed)$\times$(size). Then compute the acceleration and the sideways drift over about one second of flight.

\sol{fluids}{12} Fix the volume first: it forces the mean radius to shrink slightly, by order $\delta^2$. Then expand the area to order $\delta^2$, keeping both the change in mean radius and the slope term $\sqrt{1+r'^2}$.

\solchapter{chaos}

\sol{chaos}{10} Count dimensions of phase space. For the driven pendulum, treat the phase of the drive as a third variable. For maps, ask whether a one-dimensional map must be invertible.

\sol{chaos}{11} Subtract the two equations. You get a single equation $\dot\phi = \Delta\omega - 2K\sin\phi$ on a circle. When does it have a fixed point?

\sol{chaos}{12} Divide successive intervals. For $r_\infty$, sum the geometric series of all remaining intervals after $r_4$.

\solchapter{comp}

\sol{comp}{11} Write the step as a $2\times2$ matrix $M$. A quadratic form $\mathbf{u}^TS\mathbf{u}$ is conserved exactly when $M^TSM = S$. For the instability, look at the trace of $M$.

\sol{comp}{12} The difference between runs starts at round-off level, about $10^{-15}$, and grows like $e^{0.9t}$. For the second part, ask whether the computed trajectory might be an exact trajectory of the true system, just from a different starting point.

\solchapter{qm}

\sol{qm}{12} The bound state must decay as $e^{-\kappa\abs x}$ on both sides. The delta function produces a kink: integrate the Schr\"odinger equation from $-\epsilon$ to $+\epsilon$ to find the jump in $\psi'$.

\sol{qm}{13} Multiply the tunnelling probability by the Boltzmann factor $e^{-E/\kB T}$. One falls with energy and the other rises: find the maximum of their product.

\sol{qm}{14} Follow the amplitudes through each beam splitter. With the bomb in place, one arm acts as a measurement: the photon is either absorbed there or definitely in the other arm, and then there is nothing to interfere with.
